% Appendix to sec:method-sensitivity (content/methodology/05_statistical_model.tex): recalibration of a
%   score on the cloglog scale, SCORE2's own recalibration as a worked example, and the same formula
%   with a new outcome.
% 2026-11-16 (user: "Shrink and move"): the general text of the former Methodology subsection
%   "Recalibration of SCORE2" moved here as the first section, with its label sec:method-recalibration
%   and eq:method-recal, since the recalibrated score is only a sensitivity analysis. Its first
%   paragraph is unchanged. Of its second paragraph only the last sentence (relative weights) was
%   kept; the sentences on how SCORE2 is built and regionally recalibrated repeated the worked
%   example below. The worked example got its own section heading (app:score2-example).
%   CITATION WANTED: a general recalibration source (e.g. Janssen 2008, Steyerberg 2019); none in
%   refs.bib. CHECK: the model script enters standardised log(SCORE2), not cloglog(SCORE2)
%   (progression_who_pooled.py, design); align it before the sensitivity analysis is run.
% Written 2026-11-14 (user: "put additional information in appendix like an example of how score2
%   recalibrate itself, as we recalibrate"), from score2_loglog_explainer.tex (study note).
% Source: score2_2021, SCORE2 Updated Supplementary Material (project root, .docx, read 2026-11-14):
%   Box 1 (Fine and Gray models, stratified by cohort); Box 6 (expected vs predicted risk across age
%   groups, intercept and slope of the regression line as recalibration factors); Supplementary Methods
%   Table 2 (baseline survival 0.9605 men, 0.9776 women; centring: age 60, SBP 120, TC 6, HDL 1.3);
%   Table 3 (scale1, scale2 per region and sex); Table 4 (worked example, all numbers below).
% Recomputed in Python 2026-11-14: 1 - 0.9605^exp(0.3221) = 0.0541; ln(-ln(1 - 0.0541)) = -2.889;
%   log(-log 0.9605) + 0.3221 = -2.889; -0.5699 + 0.7476 x (-2.889) = -2.730; exp = 0.0652;
%   1 - exp(-0.0652) = 0.0631 (supplement: 0.0631).
% CAREFUL: the low-risk region is SCORE2's group for Denmark; no CGPS estimate (alpha, gamma) here,
%   those belong to Results.
\chapter{Recalibration of SCORE2}
\label{app:score2-recal}

\section{Recalibration on the Complementary Log-Log Scale}
\label{sec:method-recalibration}

A published risk score gives each person a risk $r_i$ of the outcome it was developed for, in the population it was developed in.
To predict another outcome or another population, the score can be recalibrated by taking its risk to the cloglog scale and entering it in the linear predictor with an estimated intercept $\alpha$ and slope $\delta$,
\begin{equation}
  \log\{-\log(1 - p_i)\} = \alpha + \delta \log\{-\log(1 - r_i)\}.
  \label{eq:method-recal}
\end{equation}
Fixing $\delta = 1$ and estimating only $\alpha$ adjusts the overall level of risk, referred to as calibration in the large.
Both sides of Eq.~\eqref{eq:method-recal} are log cumulative hazards, so $\delta$ states how much of the differences between persons carries over to the new outcome: a person whose score hazard is twice that of another has $2^{\delta}$ times the hazard of the new outcome, so $\delta = 1$ keeps the differences and $\delta < 1$ shrinks them.
Recalibration keeps the relative weights of the risk factors, so a new covariate correlated with a factor that the score weights wrongly for the new outcome can appear informative only by correcting that weight.

\section{Worked Example of SCORE2}
\label{app:score2-example}

SCORE2 estimates the 10-year risk of fatal and non-fatal cardiovascular disease from age, smoking, systolic blood pressure, and total and HDL cholesterol, together with the interaction of each factor with age \cite{score2_2021}.
Its coefficients were estimated with sex-specific Fine and Gray models, so its hazard is a subdistribution hazard that accounts for death from other causes, not the hazard of a model without competing events, and its risk is recalibrated to each European region on the cloglog scale \cite[Supplementary Methods, Box~1 and Box~6]{score2_2021}.
This section follows the worked example of the SCORE2 supplement, a 50-year-old man who smokes, with a systolic blood pressure of 140\,mmHg, a total cholesterol of 6.3\,mmol/L and an HDL cholesterol of 1.4\,mmol/L, living in the low-risk region \cite[Supplementary Methods, Table~4]{score2_2021}.

The centred risk factors and their coefficients sum to a linear predictor of $\mathrm{LP} = 0.3221$.
With the baseline 10-year survival of men, $S_0 = 0.9605$, the uncalibrated risk is
\begin{equation*}
  r = 1 - S_0^{\exp(\mathrm{LP})} = 1 - 0.9605^{\exp(0.3221)} = 0.0541 .
\end{equation*}
Taken to the cloglog scale, this risk is $\log\{-\log(1 - 0.0541)\} = -2.889$, which equals $\log(-\log 0.9605) + 0.3221$, the constant of the baseline survival plus the linear predictor.
The regional recalibration then applies an intercept and a slope on this scale and returns the result to a risk,
\begin{equation*}
  \text{calibrated risk} = 1 - \exp\Bigl[-\exp\bigl(\mathrm{scale}_1 + \mathrm{scale}_2 \log\{-\log(1 - r)\}\bigr)\Bigr],
\end{equation*}
with $\mathrm{scale}_1 = -0.5699$ and $\mathrm{scale}_2 = 0.7476$ for men in the low-risk region \cite[Supplementary Methods, Table~3]{score2_2021}.
Table~\ref{tab:app-score2-recal} follows the calculation step by step: the logarithm is taken twice to reach the cloglog scale, the straight line is applied there, and the exponential function is applied twice to return to a risk of 6.3\%.

\begin{table}[htbp]
  \centering
  \caption{Recalibration of the SCORE2 risk for the worked example of the SCORE2 supplement \cite[Supplementary Methods, Table~4]{score2_2021}: a 50-year-old male smoker in the low-risk region.}
  \label{tab:app-score2-recal}
  \begin{tabular}{lll}
    \toprule
    Step & Calculation & Value \\
    \midrule
    Linear predictor & sum of the centred terms & $0.3221$ \\
    Uncalibrated risk & $1 - 0.9605^{\exp(0.3221)}$ & $0.0541$ \\
    Cloglog scale & $\log\{-\log(1 - 0.0541)\}$ & $-2.889$ \\
    Intercept and slope & $-0.5699 + 0.7476 \times (-2.889)$ & $-2.730$ \\
    Cumulative hazard & $\exp(-2.730)$ & $0.0652$ \\
    Calibrated risk & $1 - \exp(-0.0652)$ & $0.0631$ \\
    \bottomrule
  \end{tabular}
\end{table}

Recalibrating SCORE2 to another outcome follows the same steps, with the intercept $\alpha$ and slope $\delta$ of Eq.~\eqref{eq:method-recal} estimated for the new outcome in place of $\mathrm{scale}_1$ and $\mathrm{scale}_2$.
The slope $\mathrm{scale}_2 = 0.7476$ of the example shows how such a slope reads: a man in the low-risk region whose uncalibrated SCORE2 hazard is twice that of another has a calibrated hazard $2^{0.7476} = 1.68$ times as high, so the recalibration narrows the differences between persons while keeping their order.
In the same way, a slope $\delta$ below 1 for a new outcome means that the score separates persons more than their risk of the new outcome differs.
